%%%%%%%%%%%%%%%%%%%%%%%%%%%%%%%%%%%%%%%%%%%%%%%%%%%%%%%%%%%%%%%%%%%%%%
% Lớp 10 | Thống kê | Mô tả và biểu diễn dữ liệu trên các bảng và biểu đồ
% Dạng 0D6?2-1: Đọc và phân tích thông tin trên bảng số liệu
% Số câu: 4  (sinh tự động bởi tach_cau.py; từ nay thầy sửa trực tiếp file này)
%%%%%%%%%%%%%%%%%%%%%%%%%%%%%%%%%%%%%%%%%%%%%%%%%%%%%%%%%%%%%%%%%%%%%%

%%%%%%%%%%%%%%%% ĐÚNG/SAI %%%%%%%%%%%%%%%%

% nguồn: MaHoaNH_DungSai_Lop10
\begin{ex}%[0D6H2-1]%[ma:85006431]
	 Điểm thi HK2 môn toán của lớp 10A9 được thông kê trong bảng dưới đây
	\begin{center}
		\begin{tabular}{|c|c|c|c|c|c|c|c|c|c|c|c|}
			\hline 
			Điểm & 4,5 & 5 & 5,5 & 6 & 6,5 & 7 & 7,5 &  8 & 8,5 & 9 & 9,5\\
			 \hline
			 Số học sinh & 2 & 3 & 3 & 4 & 5 & 5 & 6 & 4 & 3 & 3 & 2\\
			\hline 
		\end{tabular}
	\end{center}
	Các khẳng định sau  đây đúng hay sai?
	\choiceTF 
		{\True Số học sinh lớp  10A9 là $40$}
		{Tỉ lệ học sinh được trên 8 điểm  là $30\%$}
		{\True Tỉ lệ học  sinh được trên 6 điểm là $70\%$}
		{\True Tỉ lệ học sinh được ít nhất  7 điểm là $57{,}5\%$}
	 \loigiai{
		\begin{itemchoice}
			\itemch Số học sinh lớp 10A9 là $2 + 3 + 3 + 4 + 5 + 5 + 6 + 4 +  3 + 3 + 2= 40$.
			\itemch Tỉ lệ học sinh được trên 8 điểm  là $ \dfrac{8}{40} = 20\%$.
			\itemch Tỉ lệ học sinh được  trên 6 điểm là $ \dfrac{28}{40} = 70\%$.
			\itemch Tỉ lệ học sinh được ít nhất 7 điểm là  $\dfrac{23}{40} = 57{,}5\%$.
		\end{itemchoice}
	} 
\end{ex}

% nguồn: MaHoaNH_DungSai_Lop10
\begin{ex}%[0D6H2-1]%[ma:5e9baa70]
	Điểm thi HK1 môn  Văn của lớp 11B9 và 11B11 được thông kê trong bảng dưới đây
	\begin{center}
		 \textbf{Lớp 11B9} \\[8pt]
		\begin{tabular}{|c|c|c|c|c|c|c|c|c|c|c|}
			 \hline
			Điểm & 4,5 & 5 & 5,5 & 6 & 6,5 & 7 & 7,5 & 8 & 8,5  & 9\\
			 \hline
			Số học sinh & 1 & 4 & 4 & 3 & 5 & 8 & 6  & 5 & 3 & 1\\
			\hline 
		\end{tabular} \\[12pt]
		 \textbf{Lớp 11B11} \\[8pt]
		\begin{tabular}{|c|c|c|c|c|c|c|c|c|c|c|c|}
			 \hline
			Điểm & 4,5 & 5 & 5,5 & 6 & 6,5 & 7 & 7,5 & 8 & 8,5 &  9 & 9,5\\
			 \hline
			Số học sinh & 2 & 4 & 6 & 5 & 6 & 5 & 8  & 3 & 2 & 2 & 1\\
			 \hline
		\end{tabular}
	\end{center}
	Các khẳng định sau đây  đúng hay sai?
	 \choiceTF
		{\True Học sinh lớp 11B11 nhiều hơn  học sinh lớp 11B9}
		{\True Tỉ  lệ học sinh được ít nhất 7 điểm của lớp 11B9 lớn hơn lớp 11B11}
		{\True Tỉ lệ  học sinh được dưới 6 điểm của lớp 11B9 nhỏ hơn lớp 11B11}
		{Tỉ lệ học sinh được  ít nhất 8 điểm của lớp 11B9 nhỏ hơn lớp 11B11}
	\loigiai{ 
		\begin{itemchoice}
			\itemch Số học sinh lớp 11B9  là $40$, số học sinh lớp 11B11 là $44$.
			\itemch Tỉ lệ học sinh được ít nhất 7 điểm của lớp 11B9 là $\dfrac{23}{40} = 57,5 \%$.  \\
			Tỉ lệ học sinh được ít nhất 7 điểm của lớp  11B11 là $\dfrac{18}{44} \approx 47{,}7 \%$. \\
			Vậy tỉ lệ học sinh được ít nhất 7 điểm của lớp 11B9 lớn hơn lớp  11B11.
			\itemch  Tỉ lệ học sinh được dưới 6 điểm của lớp 11B9 là $\dfrac{9}{40} = 22{,}5 \%$. \\
			 Tỉ lệ học sinh được dưới 6 điểm của lớp 11B11 là $\dfrac{9}{40} = 27{,}3 \%$. \\
			Vậy tỉ  lệ học sinh được dưới 6 điểm của lớp 11B9 nhỏ hơn lớp 11B11.
			\itemch Tỉ lệ học sinh được ít nhất 8 điểm của lớp 11B9 là  $\dfrac{9}{40} = 22{,}5 \%$. \\
			Tỉ lệ học sinh được ít nhất 8  điểm của lớp 11B11 là $\dfrac{8}{44} \approx 18{,}2 \%$. \\
			Vậy tỉ lệ học sinh được ít nhất  8 điểm của lớp 11B9 lớn hơn lớp 11B11.
		\end{itemchoice}
	 }
\end{ex}


%%%%%%%%%%%%%%%% TRẢ LỜI NGẮN %%%%%%%%%%%%%%%%

% nguồn: MaHoaNH_TraLoiNgan_Lop10
\begin{bt}%[0D6N2-1]%[ma:ff2b0272]
	Trong $6$ tháng đầu năm, số sản phẩm bán ra mỗi tháng của một cửa hàng đều tăng khoảng $25\%$ so với tháng trước đó. Biết rằng, trong bảng dưới đây, số sản  phẩm bán ra của một tháng bị nhập sai. Hãy tìm tháng đó.
	\begin{center}
		\begin{tabular}{|c|c|c|c|c|c|c|}
			\hline 
			Tháng & $1$ & $2$ & $3$ & $4$ &  $5$ & $6$\\
			\hline 
			Số sản phẩm bán  ra & $145$ & $180$ & $225$ & $279$ & $390$ & $435$\\
			 \hline
		\end{tabular}
	\end{center}
	\shortans{5} 
	 \loigiai{
		Tỉ lệ phần trăm tăng thêm  của số sản phẩm bán ra mỗi tháng được tính ở bảng dưới đây
		\begin{center}
			\begin{tabular}{|c|c|c|c|c|c|}
				\hline 
				Tháng & $2$ & $3$ & $4$ & $5$  & $6$\\
				\hline 
				Tỉ lệ phần trăm tăng  thêm so với tháng trước & $24{,}1\%$ & $25\%$ & $24\%$ & $39{,}8\%$ & $11{,}5\%$\\
				\hline 
			\end{tabular}
		\end{center}
		Tỉ lệ  tăng của tháng $5$ và tháng $6$ đều rất khác so với $25\%$, do đó số liệu trong tháng $5$ là dữ liệu đã nhập sai.
	 }
\end{bt}

% nguồn: MaHoaNH_TraLoiNgan_Lop10
\begin{bt}%[0D6H2-1]%[ma:dca0e5aa]
	Trong $7$ tháng đầu năm, số sản phẩm sản xuất mỗi tháng của công ty X đều tăng trưởng khoảng $5\%$ so với tháng trước đó. Biết rằng, trong bảng dưới đây,  số sản phẩm sản xuất của một tháng bị nhập sai, Hãy tìm tháng đó.
	\begin{center}
		\begin{tabular}{|c|c|c|c|c|c|c|c|}
			\hline 
			Tháng & $1$ & $2$ & $3$ & $4$ & $5$ & $6$  & $7$\\
			 \hline
			Số sản phẩm  sản xuất & $500$ & $525$ & $551$ & $569$ & $606$ & $636$ & $668$\\
			 \hline
		\end{tabular}
	\end{center}
	 \shortans{4}
	 \loigiai{
		Tỉ lệ phần trăm tăng thêm của số sản phẩm bán ra mỗi tháng được tính ở bảng dưới  đây
		\begin{center}
			\begin{tabular}{|c|c|c|c|c|c|c|}
				 \hline
				Tháng & $2$ & $3$ & $4$ & $5$  & $6$ & $7$\\
				 \hline
				Tỉ  lệ phần trăm tăng thêm so với tháng trước & $5\%$ & $4{,}95\%$ & $3{,}4\%$ & $6{,}5\%$ & $4{,}95\%$ & $5{,}03\%$\\
				 \hline
			\end{tabular}
		\end{center}
		Ta thấy tỉ lệ  tăng của tháng $4$ và tháng $5$ đều khác xa $5\%$.\\
		Do đó trong  bảng số liệu đã cho, số sản phẩm của tháng $4$ là không chính xác.
	} 
\end{bt}
